

\begin{frame}{ElasticSearch est une serveur REST}

ElasticSearch a été conçu comme un serveur qui suit les règles
du WEb.

\begin{redbullet}

  \item Il est constitué de \alert{ressources}.

    Une ressource est une entité fournissant des services et communiquant
    par message.
    
   
   \item Les ressources sont identifiées et accessibles à des \alert{URL}.
   
   \url{https://www.example.com:9200/chemin/vers/doc?nom=b3d&type=json\#fragment}


   \item Les messages sont codés selon un \alert{protocole}, HTTP.
   

   \item Un ``message'' est une enveloppe dont le contenu est un \alert{document}.
   
   Souvent le document est en HTML,  mais cela peut être un \alert{document structuré}
   comme JSON.

\end{redbullet}

\end{frame}

\begin{frame}[fragile]
\frametitle{L'architecture REST}

REST = une définition de services Web basée sur les standards du Web.

\begin{redbullet}
  \item on s'adresse à des ressources qui fournissent des services;
  \item les messages sont échangés avec HTTP, et sont (le plus souvent) \alert{structurés} en XML ou JSON;
  \item les documents structurés fournis par une ressource sont (le plus souvent)
     produits \alert{à la volée} (par \alert{calcul}: distinction entre service et ressource statique).
\end{redbullet}

Très répandu! Excellent moyen de récupérer des informations directement exploitables
par une machine (pas comme HTML).
\end{frame}

\begin{frame}{Ressources et opérations}


\figSlideWithSize{REST}{12cm}

\vfill

Quatre opérations (celles de HTTP).
\begin{description}
  \item[\texttt{GET}] \alert{lit} (la représentation d')une ressource.
  \item[\texttt{PUT}] \alert{crée} une nouvelle ressource (ou, plus flexible: la remplace).
  \item[\texttt{POST}] envoie un message (\alert{demande de service}) à une ressource.
  \item[\texttt{DELETE}] \alert{détruit} une ressource.
\end{description}
 
\end{frame}



\begin{frame}[fragile]{Les services ElasticSearch}

ElasticSearch parle avec nous en HTTP ! C'est un service REST
qui échange des documents JSON avec nous. Essayez url{https://localhost:9200}

\vfill

Après connexion, vous devriez obtenir un beau message JSON en retour.

\vfill 
\centerline{\figSlideWithSize{elasticvue-rest.png}{9cm}}

\vfill 
Plus complet\,: les fenêtres d'ElastiVue.

\end{frame}


\begin{frame}{Application\,: la création d'un index}

On envoie un \texttt{PUT} à \url{https:/localhost:9200/films}

\vfill 

\centerline{\figSlideWithSize{es-create-index-films.png}{10cm}}

\vfill

Le serveur répond avec un message d'acquittement en JSON.

\end{frame}



\begin{frame}{Création de documents}

On envoie un \texttt{PUT} à \url{https:/localhost:9200/films/_doc/<id>}

\vfill 

\centerline{\figSlideWithSize{es-create-doc-films.png}{10cm}}

\vfill

Saurez-vous comment supprimer ce document\,?

\end{frame}


\begin{frame}{Chargeons un  jeu de données}

Récupérez \texttt{films\_esearch.json}
sur  {http://deptfod.cnam.fr/bd/tp/datasets/}. 

\vfill 

Envoyez un POST à l’URL \url{https://localhost:9200/\_bulk/} avec le contenu.
\vfill

\centerline{\figSlideWithSize{es-load-films.png}{10cm}}

\end{frame}



\begin{frame}{À vous de jouer}

\vfill 

Reproduisez les commandes d'insertion et créez la base films.

\vfill

Faites la même chose pour la base des cas.

\vfill

Amusez-vous ensuite avec ElasticVue à inspecter les index et leur contenu.

\end{frame}


